\documentclass[10pt,a4paper]{amsart}
\usepackage[margin=19mm]{geometry}
\usepackage{amsmath,amssymb,mathtools}
\usepackage[scr=rsfs,cal=euler]{mathalfa}
\usepackage{xfrac}
\usepackage[unicode,hidelinks,bookmarks=false]{hyperref}
\newcommand{\ie}{i.e.\ }
\newcommand{\cf}{cf.\ }
\newcommand{\resp}{respectively\ }
\newcommand{\Subsection}{\subsection}
\providecommand{\nobreakdash}{\nobreak\hbox{-}}
\setlength{\emergencystretch}{2em}
\multlinegap0pt
\usepackage{bm}
\usepackage{bbm}
\usepackage{dsfont}
\usepackage{xspace}
\usepackage{cancel}
\usepackage{mathtools}
\usepackage{amsthm}
\makeatletter
\@ifundefined{theorem}{\newtheorem{theorem}{Theorem}}{}
\@ifundefined{lemma}{\newtheorem{lemma}[theorem]{Lemma}}{}
\makeatother
\newcommand{\Z}{ {\mathbb{Z}} } 
\title[The size of $2$-Selmer groups for the $\frac{\pi}{3}$-congru]{The size of $2$-Selmer groups for the $\frac{\pi}{3}$-congruent number problem}
\author{Kushal Bhowmick, Aprameyo Pal}
\date{Source: arXiv:2602.08912v2 (CC BY 4.0)}
\begin{document}
\maketitle

\noindent\emph{Source and licence.} The size of $2$-Selmer groups for the $\frac{\pi}{3}$-congruent number problem. Version arXiv:2602.08912v2 (CC BY 4.0), distributed under the Creative Commons Attribution 4.0 International licence, as recorded in the arXiv licence metadata for this exact version and shown on the abstract page https://arxiv.org/abs/2602.08912v2. Full rights notice: licenses/arxiv-2602.08912-CC-BY-4.0.txt.\par\smallskip

\noindent\textit{Editorial scope.} Counted unit: the Lemma whose source label is \texttt{important lemma} in the article (source0.tex lines 293--295, source label \texttt{important lemma}), delivered together with the article's own complete original proof of it (source0.tex lines 296--326, one \texttt{proof} environment of 31 lines). No mathematics is written, repaired or re-derived: statement and proof are the article's own text line for line. Required context retained uncounted: coset (lines 285--287). Every definition, display and label of those lines is retained unchanged. Omitted from this package and not redistributed: 1--284, 288--292, 327--811. Nothing inside a retained excerpt is omitted or altered: every retained body is delivered exactly as the article's lines are written, so no omission receipt of any kind accompanies this package. Citations: the retained excerpts carry no citation command at all, so no bibliography entry of the article is transcribed and no reference is redistributed. Reference adaptation: none was needed and none was made; every reference inside the delivered unit resolves inside it, so no unresolved reference or citation is produced. Printed slips kept literal: no printed word, symbol, hypothesis or number of the counted statement or of its proof was altered. Layout and class: the article's own class is \texttt{article} with packages including authblk, inputenc, hyperref, geometry, graphicx, amssymb, amsthm, amsmath, xy, color; the wrapper is the lane's pinned \texttt{amsart} editorial wrapper at 10pt on A4, which restates the theorem-like environments the retained text needs and adds the article's own macro definitions that the retained text needs (\texttt{\textbackslash Z}), with extra wrapper packages (bm, bbm, dsfont, xspace, cancel, mathtools). The delivered numbering is the unit's own. Assets: no retained span contains a figure environment, a graphics-inclusion command, a PDF-page inclusion, a table or a TikZ picture; no image file is copied into this package, the package declares no asset dependency and no image file is redistributed with it. Licence: version arXiv:2602.08912v2 of this article is distributed under the Creative Commons Attribution 4.0 International licence, as recorded in the arXiv licence metadata for this exact version and shown on the abstract page https://arxiv.org/abs/2602.08912v2. A full rights notice is delivered at licenses/arxiv-2602.08912-CC-BY-4.0.txt. Characters: every retained excerpt is pure ASCII, so the delivered engine, XeLaTeX with its default Latin Modern text font, reports no missing character.
\begin{align}\label{coset}
&(x,y),\quad (x,y)+(0,0)=\left(\frac{-3n^{2}}{x},\frac{x+3n^{2}y}{x^{2}}\right),   \quad (x,y)+(n,0)=\left(\frac{n(x+3n)}{x-n},\frac{-4n^{2}y}{(x-n)^{2}}\right), \quad \text{and}\\&(x,y)+(-3n,0)=\left(\frac{-3n(x-n)}{x+3n},\frac{-12n^{2}y}{(x+3n)^{2}}\right) \notag.
\end{align} 

\begin{lemma}\label{important lemma}
There is a unique point $(x,y)$ in the coset (\ref{coset}) satisfying $x>0$ and $|x|_{2}\neq 1$.    
\end{lemma}

\begin{proof}
If $|x|_{2}\neq 1$, then we have $|\frac{n(x+3n)}{x-n}|_{2}=1$ and $|\frac{-3n(x-n)}{x+3n}|_{2}=1$. Therefore we get a unique point in the coset (\ref{coset}) satisfying the above condition. 

\noindent Now we assume that $|x|_{2}=1$. Note that it suffices to show $|x+3n|_{2}\neq |x-n|_{2}$. As $n$, $3n$ are odd and $x\equiv 1\pmod{2}$ we get
$$x-n=2x^{\prime}, \quad x+3n=2x^{\prime \prime} \quad \mathrm{for\: some} \quad x^{\prime},x^{\prime \prime}\in \mathbb{Z}_{2}.$$ Thus $2n=x^{\prime \prime}-x^{\prime}$. Since $n\equiv 5\pmod{8}$ we obtain
$$x^{\prime \prime}-x^{\prime}\equiv 2\pmod{8}.$$
If $x^{\prime\prime}\equiv 0,2,4,6\pmod{8}$, then we are done. So  assume that $x^{\prime\prime}\equiv 1,3,5,7\pmod{8}$. This implies that $$x\equiv 3,7\pmod{8}.$$
Below we show that this leads to a contradiction.


\noindent Now we write $\, x=\frac{r}{s},\quad y=\frac{t}{u}$ with $(r,s)=(t,u)=1$ and where $r$, $s$ are odd. Then
\begin{align*}
\frac{t^{2}}{u^{2}}  & =\frac{r}{s}(\frac{r}{s}+3n)(\frac{r}{s}-n), \\ 
\Rightarrow s^{3}t^{2} &  =r(r+3sn)(r-sn)u^{2}.
\end{align*}
Now $u^{2}\mid s^{3}$, since $t$ and $u$ are coprime. Similarly we observe that $(s^{3}, r(r+3sn)(r-sn))=1$ as $r$ and $s$ are coprime, so we have $s^{3}\mid u^{2}$. Therefore $s^{3}=u^{2}$, and it implies that $s=W^{2}$, $u=W^{3}$ for some $W \in \Z$. We now have $$r(r+3sn)(r-sn)=t^{2}.$$ 
As $r+3sn$, $r-sn$ are both even, we get $2\mid t$. So we have
\begin{equation}\label{equn}
r\left(\frac{r+3sn}{2}\right)\left(\frac{r-sn}{2}\right)=\left(\frac{t}{2}\right)^{2}.    
\end{equation}
Case 1)  $x\equiv 3\pmod{8}$. Since $s=W^{2}$, we have $s\equiv 1\pmod{8}$. Therefore $r\equiv 3\pmod{8}$ and $\frac{r-sn}{2}\equiv 3\pmod{4}$. We write $\frac{r-sn}{2}=4k+3$. Therefore $\frac{r+3sn}{2}=4k+3+2sn$. Also note that
$$\frac{r+3sn}{2}-\frac{r-sn}{2}=2sn\equiv 2\pmod{8}.$$
Putting all these together we get
$$r\left(\frac{r+3sn}{2}\right)\left(\frac{r-sn}{2}\right)=r(4k+3+2ns)(4k+3)\equiv 5\pmod{8}.$$
Using this in (\ref{equn}) we obtain a contradiction.\\

\noindent Case 2)  $x\equiv 7\pmod{8}$. Since $s=W^{2}$, we have $s\equiv 1\pmod{8}$. Therefore $r\equiv 7\pmod{8}$ and $\frac{r-sn}{2}\equiv 1\pmod{4}$. We write  $\frac{r-sn}{2}=4k+1$. Therefore $\frac{r+3sn}{2}=4k+1+2sn$. We obtain $\frac{r+3sn}{2}$, $\frac{r-sn}{2}$ are both odd. As in the previous case, we have $2sn\equiv 2\pmod{8}$. Therefore
$$r\left(\frac{r+3sn}{2}\right)\left(\frac{r-sn}{2}\right)=r(4k+1+2ns)(4k+1)\equiv 5\pmod{8}$$
which is again a contradiction.\\
\noindent Hence proved.
\end{proof}

\end{document}
